\chapter{Traçabilité et reproduction documentaire}
\label{app:reproduction}
\section{Organisation des fichiers et commandes}
Le point d'entrée est \path{report/main.tex}. Les chapitres utilisent
les sections de \path{report/sections/}, les tableaux issus de
\path{report/generated/} et les PDF déjà archivés dans
\path{results/extended/figures/}. La bibliographie est locale au
rapport. La présentation a pour point d'entrée \path{slides/main.tex}
et ses notes sont dans \path{slides/speaker_notes.md}.
Depuis la racine du projet, les commandes documentaires suivantes
n'appellent aucun solveur, runner ou calcul de métriques :
\begin{lstlisting}
.\.venv\Scripts\python.exe -B scripts/export_report_material.py
.\.venv\Scripts\python.exe -B scripts/export_report_material.py --check
.\.venv\Scripts\python.exe -B scripts/check_documentation.py
\end{lstlisting}
Le premier script lit les CSV et JSON archivés, contrôle les agrégats
et exporte les valeurs sélectionnées. Le mode \texttt{--check} compare
en mémoire tous les exports existants aux sources sans les réécrire.
Le second contrôle notamment les références, chemins et empreintes
protégées. Les résultats de ces contrôles ne valent pas compilation LaTeX.
Avec une distribution LaTeX fournissant les outils nécessaires, la
compilation peut être effectuée depuis la racine :
\begin{lstlisting}
latexmk -cd -pdf -interaction=nonstopmode -halt-on-error report/main.tex
latexmk -cd -pdf -interaction=nonstopmode -halt-on-error slides/main.tex
\end{lstlisting}
L'option \texttt{-cd} respecte les chemins relatifs depuis le répertoire
de chaque point d'entrée. Les alternatives avec \texttt{pdflatex},
\texttt{bibtex} ou Tectonic, les commandes effectivement exécutées et
les contrôles des journaux et de la mise en page sont consignés dans
\path{docs/phase5_build.md}. Ce suivi distingue la validation des
sources, la compilation et l'inspection du rendu ; leurs garanties ne
sont pas interchangeables.

Les PDF de cette édition ont été compilés et inspectés dans un
\emph{environnement externe} ; le blocage du compilateur Windows reste
consigné séparément. Les commandes, versions, vérifications et corrections
purement documentaires de cette édition figurent dans
\path{docs/phase5_v2_external_compilation.json}. Aucun solveur ni
calcul expérimental nouveau n'accompagne cette compilation.

\section{Correspondance des valeurs et limites de l'archive}
Les CSV documentaires conservent les chaînes numériques originales,
sans arrondi préalable. Le fichier
\path{report/generated/source_values.json} donne pour chaque ligne
sélectionnée la ligne de \path{summary.csv}, les valeurs sources et
les empreintes des sept fichiers qui alimentent les exports. L'arrondi
des tableaux LaTeX n'intervient qu'à l'affichage. Les paramètres $K$
sont conservés explicitement, notamment $K=0{,}025$.
Les contrôles documentaires recalculent les moyennes et écarts-types
à partir des valeurs déjà présentes dans \path{metrics.csv}, y compris
les durées enregistrées. Ils ne créent pas de mesures de solveur nouvelles.
La vérification scientifique stricte archivée reste celle de phase 4 ;
la reproduction numérique inter-environnements mentionnée dans le
rapport provient de l'audit indépendant communiqué.
L'archive de revue de phase 5 réunit les sources documentaires et les
sorties nécessaires à leur lecture. Elle exclut volontairement les
grands tableaux NPZ déjà audités. Un manifeste expérimental peut donc
y référencer un NPZ absent : cette archive documentaire n'est pas une
copie complète permettant de relancer seule le vérificateur scientifique.
Les fichiers protégés sont contrôlés dans le projet d'origine, avant
et après rédaction.

\chapter{Grille exhaustive des mesures agrégées}
\label{app:all-tables}
Les six tables suivantes reproduisent les 438 groupes du CSV agrégé.
Chaque valeur est une moyenne accompagnée de l'écart-type
échantillonnal sur dix graines. Les lignes à $n=0$ des solveurs sont
conservées telles qu'archivées : elles évaluent la même observation,
sans constituer des réalisations supplémentaires. Le bruit non traité
est également référencé à $n=0$. Les durées, MSE et erreurs relatives
complètes restent disponibles dans le CSV associé à cet export.
% All 438 archived summary groups, including n=0; sample SD, ten seeds.
\begin{longtable}{llrll}
\caption{Régions géométriques, $\sigma=0.025$ : grille complète, moyennes $\pm$ écarts-types sur dix graines.}\\
\toprule
Méthode & $K$ & $n$ & PSNR (dB) & SSIM \\
\midrule
\endfirsthead
\toprule
Méthode & $K$ & $n$ & PSNR (dB) & SSIM \\
\midrule
\endhead
\bottomrule
\endfoot
Chaleur & -- & 0 & $32.0401 \pm 0.0458$ & $0.6646 \pm 0.0020$ \\
Chaleur & -- & 1 & $30.4691 \pm 0.0216$ & $0.8740 \pm 0.0015$ \\
Chaleur & -- & 2 & $29.2461 \pm 0.0190$ & $0.9074 \pm 0.0012$ \\
Chaleur & -- & 5 & $26.8719 \pm 0.0150$ & $0.9007 \pm 0.0008$ \\
Chaleur & -- & 10 & $25.2946 \pm 0.0124$ & $0.8708 \pm 0.0006$ \\
Chaleur & -- & 20 & $23.7481 \pm 0.0101$ & $0.8245 \pm 0.0005$ \\
Chaleur & -- & 40 & $22.2170 \pm 0.0088$ & $0.7696 \pm 0.0003$ \\
Chaleur & -- & 80 & $20.6797 \pm 0.0082$ & $0.7159 \pm 0.0002$ \\
Bruitée & -- & 0 & $32.0401 \pm 0.0458$ & $0.6646 \pm 0.0020$ \\
PM exp. & 0.025 & 0 & $32.0401 \pm 0.0458$ & $0.6646 \pm 0.0020$ \\
PM exp. & 0.025 & 1 & $32.6135 \pm 0.0536$ & $0.6923 \pm 0.0023$ \\
PM exp. & 0.025 & 2 & $33.1924 \pm 0.0623$ & $0.7195 \pm 0.0026$ \\
PM exp. & 0.025 & 5 & $34.9358 \pm 0.0919$ & $0.7938 \pm 0.0033$ \\
PM exp. & 0.025 & 10 & $37.5746 \pm 0.1342$ & $0.8795 \pm 0.0037$ \\
PM exp. & 0.025 & 20 & $40.5494 \pm 0.2173$ & $0.9374 \pm 0.0035$ \\
PM exp. & 0.025 & 40 & $42.6393 \pm 0.3200$ & $0.9605 \pm 0.0033$ \\
PM exp. & 0.025 & 80 & $44.2669 \pm 0.3844$ & $0.9724 \pm 0.0029$ \\
PM exp. & 0.05 & 0 & $32.0401 \pm 0.0458$ & $0.6646 \pm 0.0020$ \\
PM exp. & 0.05 & 1 & $34.4836 \pm 0.0672$ & $0.7741 \pm 0.0025$ \\
PM exp. & 0.05 & 2 & $37.1016 \pm 0.0932$ & $0.8643 \pm 0.0027$ \\
PM exp. & 0.05 & 5 & $43.3047 \pm 0.1650$ & $0.9698 \pm 0.0013$ \\
PM exp. & 0.05 & 10 & $47.5433 \pm 0.2525$ & $0.9916 \pm 0.0005$ \\
PM exp. & 0.05 & 20 & $50.7729 \pm 0.4245$ & $0.9974 \pm 0.0002$ \\
PM exp. & 0.05 & 40 & $53.5923 \pm 0.6869$ & $0.9992 \pm 0.0001$ \\
PM exp. & 0.05 & 80 & $56.2055 \pm 1.0724$ & $0.9998 \pm 4.10\times10^{-5}$ \\
PM exp. & 0.1 & 0 & $32.0401 \pm 0.0458$ & $0.6646 \pm 0.0020$ \\
PM exp. & 0.1 & 1 & $37.4157 \pm 0.0737$ & $0.8720 \pm 0.0017$ \\
PM exp. & 0.1 & 2 & $41.0148 \pm 0.1195$ & $0.9445 \pm 0.0012$ \\
PM exp. & 0.1 & 5 & $45.0444 \pm 0.2204$ & $0.9825 \pm 0.0008$ \\
PM exp. & 0.1 & 10 & $46.2944 \pm 0.4012$ & $0.9916 \pm 0.0006$ \\
PM exp. & 0.1 & 20 & $42.0578 \pm 0.5082$ & $0.9865 \pm 0.0014$ \\
PM exp. & 0.1 & 40 & $37.8196 \pm 0.1081$ & $0.9688 \pm 0.0007$ \\
PM exp. & 0.1 & 80 & $35.7265 \pm 0.0441$ & $0.9576 \pm 0.0002$ \\
PM exp. & 0.2 & 0 & $32.0401 \pm 0.0458$ & $0.6646 \pm 0.0020$ \\
PM exp. & 0.2 & 1 & $38.5518 \pm 0.0812$ & $0.9012 \pm 0.0014$ \\
PM exp. & 0.2 & 2 & $40.9490 \pm 0.1235$ & $0.9504 \pm 0.0012$ \\
PM exp. & 0.2 & 5 & $41.0957 \pm 0.1190$ & $0.9742 \pm 0.0008$ \\
PM exp. & 0.2 & 10 & $39.8047 \pm 0.0816$ & $0.9757 \pm 0.0005$ \\
PM exp. & 0.2 & 20 & $38.3459 \pm 0.0525$ & $0.9707 \pm 0.0003$ \\
PM exp. & 0.2 & 40 & $36.8281 \pm 0.0331$ & $0.9628 \pm 0.0002$ \\
PM exp. & 0.2 & 80 & $35.2997 \pm 0.0275$ & $0.9557 \pm 0.0001$ \\
PM rat. & 0.025 & 0 & $32.0401 \pm 0.0458$ & $0.6646 \pm 0.0020$ \\
PM rat. & 0.025 & 1 & $33.6221 \pm 0.0571$ & $0.7376 \pm 0.0023$ \\
PM rat. & 0.025 & 2 & $35.3384 \pm 0.0705$ & $0.8066 \pm 0.0024$ \\
PM rat. & 0.025 & 5 & $40.7641 \pm 0.1321$ & $0.9414 \pm 0.0017$ \\
PM rat. & 0.025 & 10 & $46.4818 \pm 0.2400$ & $0.9885 \pm 0.0006$ \\
PM rat. & 0.025 & 20 & $49.8478 \pm 0.3996$ & $0.9968 \pm 0.0002$ \\
PM rat. & 0.025 & 40 & $50.7053 \pm 0.5756$ & $0.9984 \pm 0.0002$ \\
PM rat. & 0.025 & 80 & $48.2234 \pm 0.4522$ & $0.9978 \pm 0.0002$ \\
PM rat. & 0.05 & 0 & $32.0401 \pm 0.0458$ & $0.6646 \pm 0.0020$ \\
PM rat. & 0.05 & 1 & $35.6142 \pm 0.0656$ & $0.8160 \pm 0.0021$ \\
PM rat. & 0.05 & 2 & $39.0400 \pm 0.0945$ & $0.9110 \pm 0.0016$ \\
PM rat. & 0.05 & 5 & $44.3851 \pm 0.1986$ & $0.9789 \pm 0.0009$ \\
PM rat. & 0.05 & 10 & $46.3019 \pm 0.3011$ & $0.9914 \pm 0.0005$ \\
PM rat. & 0.05 & 20 & $44.7342 \pm 0.3488$ & $0.9927 \pm 0.0005$ \\
PM rat. & 0.05 & 40 & $39.7774 \pm 0.2397$ & $0.9809 \pm 0.0009$ \\
PM rat. & 0.05 & 80 & $35.6583 \pm 0.0714$ & $0.9586 \pm 0.0003$ \\
PM rat. & 0.1 & 0 & $32.0401 \pm 0.0458$ & $0.6646 \pm 0.0020$ \\
PM rat. & 0.1 & 1 & $37.7415 \pm 0.0753$ & $0.8813 \pm 0.0015$ \\
PM rat. & 0.1 & 2 & $40.9177 \pm 0.1225$ & $0.9464 \pm 0.0012$ \\
PM rat. & 0.1 & 5 & $42.4115 \pm 0.1872$ & $0.9783 \pm 0.0009$ \\
PM rat. & 0.1 & 10 & $39.9642 \pm 0.1549$ & $0.9789 \pm 0.0006$ \\
PM rat. & 0.1 & 20 & $36.8146 \pm 0.0933$ & $0.9664 \pm 0.0005$ \\
PM rat. & 0.1 & 40 & $33.4343 \pm 0.0713$ & $0.9439 \pm 0.0005$ \\
PM rat. & 0.1 & 80 & $29.0777 \pm 0.0604$ & $0.9025 \pm 0.0006$ \\
PM rat. & 0.2 & 0 & $32.0401 \pm 0.0458$ & $0.6646 \pm 0.0020$ \\
PM rat. & 0.2 & 1 & $38.2708 \pm 0.0843$ & $0.9008 \pm 0.0014$ \\
PM rat. & 0.2 & 2 & $39.5296 \pm 0.1097$ & $0.9477 \pm 0.0012$ \\
PM rat. & 0.2 & 5 & $36.7299 \pm 0.0926$ & $0.9638 \pm 0.0008$ \\
PM rat. & 0.2 & 10 & $32.4654 \pm 0.0675$ & $0.9471 \pm 0.0006$ \\
PM rat. & 0.2 & 20 & $27.2821 \pm 0.0391$ & $0.8938 \pm 0.0007$ \\
PM rat. & 0.2 & 40 & $23.4187 \pm 0.0152$ & $0.8085 \pm 0.0005$ \\
PM rat. & 0.2 & 80 & $21.2150 \pm 0.0096$ & $0.7336 \pm 0.0003$ \\
\end{longtable}

\begin{longtable}{llrll}
\caption{Régions géométriques, $\sigma=0.05$ : grille complète, moyennes $\pm$ écarts-types sur dix graines.}\\
\toprule
Méthode & $K$ & $n$ & PSNR (dB) & SSIM \\
\midrule
\endfirsthead
\toprule
Méthode & $K$ & $n$ & PSNR (dB) & SSIM \\
\midrule
\endhead
\bottomrule
\endfoot
Chaleur & -- & 0 & $26.0195 \pm 0.0458$ & $0.3914 \pm 0.0017$ \\
Chaleur & -- & 1 & $28.9376 \pm 0.0428$ & $0.6973 \pm 0.0033$ \\
Chaleur & -- & 2 & $28.6107 \pm 0.0385$ & $0.8053 \pm 0.0033$ \\
Chaleur & -- & 5 & $26.7195 \pm 0.0308$ & $0.8636 \pm 0.0025$ \\
Chaleur & -- & 10 & $25.2406 \pm 0.0254$ & $0.8564 \pm 0.0015$ \\
Chaleur & -- & 20 & $23.7283 \pm 0.0208$ & $0.8197 \pm 0.0010$ \\
Chaleur & -- & 40 & $22.2094 \pm 0.0179$ & $0.7681 \pm 0.0006$ \\
Chaleur & -- & 80 & $20.6768 \pm 0.0165$ & $0.7154 \pm 0.0005$ \\
Bruitée & -- & 0 & $26.0195 \pm 0.0458$ & $0.3914 \pm 0.0017$ \\
PM exp. & 0.025 & 0 & $26.0195 \pm 0.0458$ & $0.3914 \pm 0.0017$ \\
PM exp. & 0.025 & 1 & $26.1093 \pm 0.0471$ & $0.3953 \pm 0.0017$ \\
PM exp. & 0.025 & 2 & $26.1931 \pm 0.0482$ & $0.3990 \pm 0.0018$ \\
PM exp. & 0.025 & 5 & $26.4136 \pm 0.0509$ & $0.4090 \pm 0.0019$ \\
PM exp. & 0.025 & 10 & $26.7150 \pm 0.0545$ & $0.4232 \pm 0.0021$ \\
PM exp. & 0.025 & 20 & $27.1885 \pm 0.0607$ & $0.4469 \pm 0.0024$ \\
PM exp. & 0.025 & 40 & $27.8988 \pm 0.0768$ & $0.4852 \pm 0.0034$ \\
PM exp. & 0.025 & 80 & $28.7598 \pm 0.1009$ & $0.5327 \pm 0.0056$ \\
PM exp. & 0.05 & 0 & $26.0195 \pm 0.0458$ & $0.3914 \pm 0.0017$ \\
PM exp. & 0.05 & 1 & $26.5919 \pm 0.0535$ & $0.4169 \pm 0.0021$ \\
PM exp. & 0.05 & 2 & $27.1695 \pm 0.0621$ & $0.4443 \pm 0.0026$ \\
PM exp. & 0.05 & 5 & $28.9068 \pm 0.0912$ & $0.5359 \pm 0.0043$ \\
PM exp. & 0.05 & 10 & $31.5266 \pm 0.1341$ & $0.6831 \pm 0.0070$ \\
PM exp. & 0.05 & 20 & $34.4402 \pm 0.2247$ & $0.8190 \pm 0.0085$ \\
PM exp. & 0.05 & 40 & $36.4299 \pm 0.3274$ & $0.8833 \pm 0.0087$ \\
PM exp. & 0.05 & 80 & $37.9419 \pm 0.3772$ & $0.9180 \pm 0.0082$ \\
PM exp. & 0.1 & 0 & $26.0195 \pm 0.0458$ & $0.3914 \pm 0.0017$ \\
PM exp. & 0.1 & 1 & $28.4571 \pm 0.0673$ & $0.5063 \pm 0.0030$ \\
PM exp. & 0.1 & 2 & $31.0611 \pm 0.0942$ & $0.6439 \pm 0.0046$ \\
PM exp. & 0.1 & 5 & $37.1361 \pm 0.1860$ & $0.8936 \pm 0.0042$ \\
PM exp. & 0.1 & 10 & $40.7063 \pm 0.3436$ & $0.9650 \pm 0.0019$ \\
PM exp. & 0.1 & 20 & $40.3634 \pm 0.6243$ & $0.9779 \pm 0.0026$ \\
PM exp. & 0.1 & 40 & $37.5412 \pm 0.2336$ & $0.9668 \pm 0.0015$ \\
PM exp. & 0.1 & 80 & $35.6318 \pm 0.0909$ & $0.9570 \pm 0.0005$ \\
PM exp. & 0.2 & 0 & $26.0195 \pm 0.0458$ & $0.3914 \pm 0.0017$ \\
PM exp. & 0.2 & 1 & $31.3890 \pm 0.0753$ & $0.6542 \pm 0.0032$ \\
PM exp. & 0.2 & 2 & $34.8898 \pm 0.1203$ & $0.8169 \pm 0.0033$ \\
PM exp. & 0.2 & 5 & $37.9130 \pm 0.1741$ & $0.9283 \pm 0.0028$ \\
PM exp. & 0.2 & 10 & $38.3505 \pm 0.1530$ & $0.9579 \pm 0.0015$ \\
PM exp. & 0.2 & 20 & $37.7540 \pm 0.1134$ & $0.9645 \pm 0.0008$ \\
PM exp. & 0.2 & 40 & $36.6014 \pm 0.0716$ & $0.9608 \pm 0.0005$ \\
PM exp. & 0.2 & 80 & $35.2197 \pm 0.0534$ & $0.9551 \pm 0.0003$ \\
PM rat. & 0.025 & 0 & $26.0195 \pm 0.0458$ & $0.3914 \pm 0.0017$ \\
PM rat. & 0.025 & 1 & $26.5759 \pm 0.0505$ & $0.4155 \pm 0.0019$ \\
PM rat. & 0.025 & 2 & $27.1542 \pm 0.0556$ & $0.4421 \pm 0.0022$ \\
PM rat. & 0.025 & 5 & $29.0305 \pm 0.0738$ & $0.5364 \pm 0.0034$ \\
PM rat. & 0.025 & 10 & $32.6042 \pm 0.1173$ & $0.7239 \pm 0.0053$ \\
PM rat. & 0.025 & 20 & $39.6339 \pm 0.2601$ & $0.9463 \pm 0.0034$ \\
PM rat. & 0.025 & 40 & $45.4018 \pm 0.4980$ & $0.9936 \pm 0.0005$ \\
PM rat. & 0.025 & 80 & $46.1862 \pm 0.6477$ & $0.9967 \pm 0.0004$ \\
PM rat. & 0.05 & 0 & $26.0195 \pm 0.0458$ & $0.3914 \pm 0.0017$ \\
PM rat. & 0.05 & 1 & $27.5998 \pm 0.0572$ & $0.4627 \pm 0.0024$ \\
PM rat. & 0.05 & 2 & $29.3122 \pm 0.0709$ & $0.5494 \pm 0.0032$ \\
PM rat. & 0.05 & 5 & $34.6930 \pm 0.1361$ & $0.8118 \pm 0.0045$ \\
PM rat. & 0.05 & 10 & $40.0986 \pm 0.2567$ & $0.9551 \pm 0.0021$ \\
PM rat. & 0.05 & 20 & $41.9526 \pm 0.4437$ & $0.9832 \pm 0.0012$ \\
PM rat. & 0.05 & 40 & $39.1773 \pm 0.4279$ & $0.9779 \pm 0.0019$ \\
PM rat. & 0.05 & 80 & $35.5528 \pm 0.1449$ & $0.9579 \pm 0.0007$ \\
PM rat. & 0.1 & 0 & $26.0195 \pm 0.0458$ & $0.3914 \pm 0.0017$ \\
PM rat. & 0.1 & 1 & $29.5885 \pm 0.0665$ & $0.5619 \pm 0.0030$ \\
PM rat. & 0.1 & 2 & $32.9822 \pm 0.0972$ & $0.7351 \pm 0.0035$ \\
PM rat. & 0.1 & 5 & $37.8375 \pm 0.1969$ & $0.9200 \pm 0.0030$ \\
PM rat. & 0.1 & 10 & $38.3228 \pm 0.2388$ & $0.9587 \pm 0.0017$ \\
PM rat. & 0.1 & 20 & $36.3626 \pm 0.1810$ & $0.9600 \pm 0.0011$ \\
PM rat. & 0.1 & 40 & $33.3179 \pm 0.1437$ & $0.9419 \pm 0.0010$ \\
PM rat. & 0.1 & 80 & $29.0486 \pm 0.1215$ & $0.9018 \pm 0.0011$ \\
PM rat. & 0.2 & 0 & $26.0195 \pm 0.0458$ & $0.3914 \pm 0.0017$ \\
PM rat. & 0.2 & 1 & $31.6864 \pm 0.0777$ & $0.6713 \pm 0.0029$ \\
PM rat. & 0.2 & 2 & $34.6475 \pm 0.1195$ & $0.8222 \pm 0.0032$ \\
PM rat. & 0.2 & 5 & $35.2280 \pm 0.1500$ & $0.9197 \pm 0.0026$ \\
PM rat. & 0.2 & 10 & $32.1179 \pm 0.1282$ & $0.9303 \pm 0.0016$ \\
PM rat. & 0.2 & 20 & $27.2270 \pm 0.0781$ & $0.8884 \pm 0.0015$ \\
PM rat. & 0.2 & 40 & $23.4097 \pm 0.0308$ & $0.8069 \pm 0.0010$ \\
PM rat. & 0.2 & 80 & $21.2120 \pm 0.0194$ & $0.7331 \pm 0.0005$ \\
\end{longtable}

\begin{longtable}{llrll}
\caption{Régions géométriques, $\sigma=0.1$ : grille complète, moyennes $\pm$ écarts-types sur dix graines.}\\
\toprule
Méthode & $K$ & $n$ & PSNR (dB) & SSIM \\
\midrule
\endfirsthead
\toprule
Méthode & $K$ & $n$ & PSNR (dB) & SSIM \\
\midrule
\endhead
\bottomrule
\endfoot
Chaleur & -- & 0 & $19.9989 \pm 0.0458$ & $0.2168 \pm 0.0013$ \\
Chaleur & -- & 1 & $25.5341 \pm 0.0665$ & $0.4306 \pm 0.0035$ \\
Chaleur & -- & 2 & $26.7244 \pm 0.0724$ & $0.5828 \pm 0.0054$ \\
Chaleur & -- & 5 & $26.1523 \pm 0.0636$ & $0.7506 \pm 0.0064$ \\
Chaleur & -- & 10 & $25.0260 \pm 0.0527$ & $0.8058 \pm 0.0041$ \\
Chaleur & -- & 20 & $23.6482 \pm 0.0432$ & $0.8017 \pm 0.0023$ \\
Chaleur & -- & 40 & $22.1790 \pm 0.0371$ & $0.7625 \pm 0.0013$ \\
Chaleur & -- & 80 & $20.6655 \pm 0.0337$ & $0.7138 \pm 0.0009$ \\
Bruitée & -- & 0 & $19.9989 \pm 0.0458$ & $0.2168 \pm 0.0013$ \\
PM exp. & 0.025 & 0 & $19.9989 \pm 0.0458$ & $0.2168 \pm 0.0013$ \\
PM exp. & 0.025 & 1 & $20.0109 \pm 0.0460$ & $0.2171 \pm 0.0013$ \\
PM exp. & 0.025 & 2 & $20.0219 \pm 0.0461$ & $0.2173 \pm 0.0013$ \\
PM exp. & 0.025 & 5 & $20.0493 \pm 0.0467$ & $0.2179 \pm 0.0013$ \\
PM exp. & 0.025 & 10 & $20.0827 \pm 0.0473$ & $0.2186 \pm 0.0013$ \\
PM exp. & 0.025 & 20 & $20.1269 \pm 0.0483$ & $0.2196 \pm 0.0013$ \\
PM exp. & 0.025 & 40 & $20.1816 \pm 0.0492$ & $0.2208 \pm 0.0013$ \\
PM exp. & 0.025 & 80 & $20.2466 \pm 0.0507$ & $0.2223 \pm 0.0013$ \\
PM exp. & 0.05 & 0 & $19.9989 \pm 0.0458$ & $0.2168 \pm 0.0013$ \\
PM exp. & 0.05 & 1 & $20.0888 \pm 0.0471$ & $0.2188 \pm 0.0013$ \\
PM exp. & 0.05 & 2 & $20.1725 \pm 0.0482$ & $0.2206 \pm 0.0013$ \\
PM exp. & 0.05 & 5 & $20.3928 \pm 0.0511$ & $0.2257 \pm 0.0014$ \\
PM exp. & 0.05 & 10 & $20.6935 \pm 0.0549$ & $0.2330 \pm 0.0015$ \\
PM exp. & 0.05 & 20 & $21.1646 \pm 0.0617$ & $0.2455 \pm 0.0018$ \\
PM exp. & 0.05 & 40 & $21.8685 \pm 0.0785$ & $0.2675 \pm 0.0025$ \\
PM exp. & 0.05 & 80 & $22.7164 \pm 0.0988$ & $0.2985 \pm 0.0041$ \\
PM exp. & 0.1 & 0 & $19.9989 \pm 0.0458$ & $0.2168 \pm 0.0013$ \\
PM exp. & 0.1 & 1 & $20.5719 \pm 0.0537$ & $0.2298 \pm 0.0015$ \\
PM exp. & 0.1 & 2 & $21.1497 \pm 0.0625$ & $0.2443 \pm 0.0018$ \\
PM exp. & 0.1 & 5 & $22.8854 \pm 0.0927$ & $0.2997 \pm 0.0032$ \\
PM exp. & 0.1 & 10 & $25.4934 \pm 0.1357$ & $0.4290 \pm 0.0075$ \\
PM exp. & 0.1 & 20 & $28.3449 \pm 0.2351$ & $0.6254 \pm 0.0137$ \\
PM exp. & 0.1 & 40 & $29.9811 \pm 0.3203$ & $0.7432 \pm 0.0162$ \\
PM exp. & 0.1 & 80 & $30.7318 \pm 0.3111$ & $0.8051 \pm 0.0156$ \\
PM exp. & 0.2 & 0 & $19.9989 \pm 0.0458$ & $0.2168 \pm 0.0013$ \\
PM exp. & 0.2 & 1 & $22.4396 \pm 0.0678$ & $0.2795 \pm 0.0021$ \\
PM exp. & 0.2 & 2 & $25.0458 \pm 0.0954$ & $0.3783 \pm 0.0040$ \\
PM exp. & 0.2 & 5 & $31.0611 \pm 0.1795$ & $0.7061 \pm 0.0084$ \\
PM exp. & 0.2 & 10 & $34.4372 \pm 0.2384$ & $0.8770 \pm 0.0046$ \\
PM exp. & 0.2 & 20 & $35.7056 \pm 0.3280$ & $0.9378 \pm 0.0026$ \\
PM exp. & 0.2 & 40 & $35.6393 \pm 0.2811$ & $0.9524 \pm 0.0013$ \\
PM exp. & 0.2 & 80 & $34.7866 \pm 0.2153$ & $0.9525 \pm 0.0008$ \\
PM rat. & 0.025 & 0 & $19.9989 \pm 0.0458$ & $0.2168 \pm 0.0013$ \\
PM rat. & 0.025 & 1 & $20.1685 \pm 0.0474$ & $0.2204 \pm 0.0013$ \\
PM rat. & 0.025 & 2 & $20.3392 \pm 0.0491$ & $0.2242 \pm 0.0014$ \\
PM rat. & 0.025 & 5 & $20.8599 \pm 0.0542$ & $0.2363 \pm 0.0015$ \\
PM rat. & 0.025 & 10 & $21.7651 \pm 0.0640$ & $0.2600 \pm 0.0019$ \\
PM rat. & 0.025 & 20 & $23.7556 \pm 0.0893$ & $0.3276 \pm 0.0032$ \\
PM rat. & 0.025 & 40 & $28.4102 \pm 0.1684$ & $0.5786 \pm 0.0097$ \\
PM rat. & 0.025 & 80 & $36.4457 \pm 0.5449$ & $0.9315 \pm 0.0075$ \\
PM rat. & 0.05 & 0 & $19.9989 \pm 0.0458$ & $0.2168 \pm 0.0013$ \\
PM rat. & 0.05 & 1 & $20.5558 \pm 0.0506$ & $0.2290 \pm 0.0014$ \\
PM rat. & 0.05 & 2 & $21.1345 \pm 0.0559$ & $0.2429 \pm 0.0016$ \\
PM rat. & 0.05 & 5 & $23.0105 \pm 0.0751$ & $0.2983 \pm 0.0025$ \\
PM rat. & 0.05 & 10 & $26.5705 \pm 0.1211$ & $0.4587 \pm 0.0059$ \\
PM rat. & 0.05 & 20 & $33.3624 \pm 0.2979$ & $0.8348 \pm 0.0090$ \\
PM rat. & 0.05 & 40 & $37.0422 \pm 0.5734$ & $0.9611 \pm 0.0035$ \\
PM rat. & 0.05 & 80 & $35.1135 \pm 0.3158$ & $0.9547 \pm 0.0016$ \\
PM rat. & 0.1 & 0 & $19.9989 \pm 0.0458$ & $0.2168 \pm 0.0013$ \\
PM rat. & 0.1 & 1 & $21.5805 \pm 0.0576$ & $0.2541 \pm 0.0017$ \\
PM rat. & 0.1 & 2 & $23.2927 \pm 0.0720$ & $0.3063 \pm 0.0024$ \\
PM rat. & 0.1 & 5 & $28.6326 \pm 0.1392$ & $0.5643 \pm 0.0067$ \\
PM rat. & 0.1 & 10 & $33.6243 \pm 0.2427$ & $0.8465 \pm 0.0057$ \\
PM rat. & 0.1 & 20 & $34.6530 \pm 0.2974$ & $0.9292 \pm 0.0027$ \\
PM rat. & 0.1 & 40 & $32.7916 \pm 0.2834$ & $0.9331 \pm 0.0022$ \\
PM rat. & 0.1 & 80 & $28.8937 \pm 0.2417$ & $0.8988 \pm 0.0023$ \\
PM rat. & 0.2 & 0 & $19.9989 \pm 0.0458$ & $0.2168 \pm 0.0013$ \\
PM rat. & 0.2 & 1 & $23.5666 \pm 0.0677$ & $0.3141 \pm 0.0023$ \\
PM rat. & 0.2 & 2 & $26.9275 \pm 0.0986$ & $0.4623 \pm 0.0041$ \\
PM rat. & 0.2 & 5 & $31.3253 \pm 0.1804$ & $0.7580 \pm 0.0066$ \\
PM rat. & 0.2 & 10 & $30.8239 \pm 0.2099$ & $0.8641 \pm 0.0044$ \\
PM rat. & 0.2 & 20 & $26.9975 \pm 0.1527$ & $0.8667 \pm 0.0033$ \\
PM rat. & 0.2 & 40 & $23.3714 \pm 0.0630$ & $0.8009 \pm 0.0022$ \\
PM rat. & 0.2 & 80 & $21.1999 \pm 0.0396$ & $0.7315 \pm 0.0010$ \\
\end{longtable}

\begin{longtable}{llrll}
\caption{Rampe et contours, $\sigma=0.025$ : grille complète, moyennes $\pm$ écarts-types sur dix graines.}\\
\toprule
Méthode & $K$ & $n$ & PSNR (dB) & SSIM \\
\midrule
\endfirsthead
\toprule
Méthode & $K$ & $n$ & PSNR (dB) & SSIM \\
\midrule
\endhead
\bottomrule
\endfoot
Chaleur & -- & 0 & $32.0401 \pm 0.0458$ & $0.6970 \pm 0.0019$ \\
Chaleur & -- & 1 & $29.5660 \pm 0.0198$ & $0.8867 \pm 0.0017$ \\
Chaleur & -- & 2 & $29.4304 \pm 0.0134$ & $0.9223 \pm 0.0013$ \\
Chaleur & -- & 5 & $27.2697 \pm 0.0102$ & $0.9180 \pm 0.0007$ \\
Chaleur & -- & 10 & $25.6750 \pm 0.0094$ & $0.8917 \pm 0.0003$ \\
Chaleur & -- & 20 & $23.9598 \pm 0.0081$ & $0.8467 \pm 0.0002$ \\
Chaleur & -- & 40 & $22.3831 \pm 0.0070$ & $0.7933 \pm 0.0003$ \\
Chaleur & -- & 80 & $21.0512 \pm 0.0057$ & $0.7488 \pm 0.0003$ \\
Bruitée & -- & 0 & $32.0401 \pm 0.0458$ & $0.6970 \pm 0.0019$ \\
PM exp. & 0.025 & 0 & $32.0401 \pm 0.0458$ & $0.6970 \pm 0.0019$ \\
PM exp. & 0.025 & 1 & $32.6052 \pm 0.0531$ & $0.7229 \pm 0.0022$ \\
PM exp. & 0.025 & 2 & $33.1729 \pm 0.0609$ & $0.7480 \pm 0.0024$ \\
PM exp. & 0.025 & 5 & $34.8617 \pm 0.0868$ & $0.8153 \pm 0.0030$ \\
PM exp. & 0.025 & 10 & $37.3861 \pm 0.1271$ & $0.8912 \pm 0.0033$ \\
PM exp. & 0.025 & 20 & $40.3640 \pm 0.2060$ & $0.9448 \pm 0.0031$ \\
PM exp. & 0.025 & 40 & $42.3182 \pm 0.3074$ & $0.9657 \pm 0.0030$ \\
PM exp. & 0.025 & 80 & $42.9393 \pm 0.3156$ & $0.9749 \pm 0.0021$ \\
PM exp. & 0.05 & 0 & $32.0401 \pm 0.0458$ & $0.6970 \pm 0.0019$ \\
PM exp. & 0.05 & 1 & $34.4466 \pm 0.0665$ & $0.7985 \pm 0.0024$ \\
PM exp. & 0.05 & 2 & $36.9977 \pm 0.0887$ & $0.8796 \pm 0.0025$ \\
PM exp. & 0.05 & 5 & $42.9282 \pm 0.1113$ & $0.9726 \pm 0.0011$ \\
PM exp. & 0.05 & 10 & $46.5380 \pm 0.1719$ & $0.9913 \pm 0.0005$ \\
PM exp. & 0.05 & 20 & $47.6194 \pm 0.2821$ & $0.9948 \pm 0.0004$ \\
PM exp. & 0.05 & 40 & $46.2645 \pm 0.2454$ & $0.9936 \pm 0.0004$ \\
PM exp. & 0.05 & 80 & $43.5241 \pm 0.1865$ & $0.9896 \pm 0.0004$ \\
PM exp. & 0.1 & 0 & $32.0401 \pm 0.0458$ & $0.6970 \pm 0.0019$ \\
PM exp. & 0.1 & 1 & $37.3081 \pm 0.0650$ & $0.8870 \pm 0.0018$ \\
PM exp. & 0.1 & 2 & $40.6919 \pm 0.0818$ & $0.9506 \pm 0.0012$ \\
PM exp. & 0.1 & 5 & $43.5262 \pm 0.1824$ & $0.9818 \pm 0.0009$ \\
PM exp. & 0.1 & 10 & $42.5843 \pm 0.2174$ & $0.9858 \pm 0.0007$ \\
PM exp. & 0.1 & 20 & $38.9928 \pm 0.1873$ & $0.9762 \pm 0.0008$ \\
PM exp. & 0.1 & 40 & $36.1982 \pm 0.1202$ & $0.9606 \pm 0.0007$ \\
PM exp. & 0.1 & 80 & $33.9862 \pm 0.1066$ & $0.9454 \pm 0.0008$ \\
PM exp. & 0.2 & 0 & $32.0401 \pm 0.0458$ & $0.6970 \pm 0.0019$ \\
PM exp. & 0.2 & 1 & $38.2183 \pm 0.0666$ & $0.9122 \pm 0.0015$ \\
PM exp. & 0.2 & 2 & $39.8545 \pm 0.1233$ & $0.9540 \pm 0.0013$ \\
PM exp. & 0.2 & 5 & $38.3901 \pm 0.1350$ & $0.9707 \pm 0.0008$ \\
PM exp. & 0.2 & 10 & $35.9755 \pm 0.1243$ & $0.9654 \pm 0.0005$ \\
PM exp. & 0.2 & 20 & $33.6242 \pm 0.0806$ & $0.9497 \pm 0.0004$ \\
PM exp. & 0.2 & 40 & $31.4394 \pm 0.0652$ & $0.9287 \pm 0.0004$ \\
PM exp. & 0.2 & 80 & $28.5844 \pm 0.1368$ & $0.8965 \pm 0.0022$ \\
PM rat. & 0.025 & 0 & $32.0401 \pm 0.0458$ & $0.6970 \pm 0.0019$ \\
PM rat. & 0.025 & 1 & $33.6013 \pm 0.0564$ & $0.7651 \pm 0.0022$ \\
PM rat. & 0.025 & 2 & $35.2827 \pm 0.0677$ & $0.8280 \pm 0.0023$ \\
PM rat. & 0.025 & 5 & $40.4722 \pm 0.0984$ & $0.9472 \pm 0.0016$ \\
PM rat. & 0.025 & 10 & $45.5535 \pm 0.1473$ & $0.9885 \pm 0.0006$ \\
PM rat. & 0.025 & 20 & $46.9468 \pm 0.2176$ & $0.9946 \pm 0.0003$ \\
PM rat. & 0.025 & 40 & $44.3949 \pm 0.1909$ & $0.9919 \pm 0.0003$ \\
PM rat. & 0.025 & 80 & $40.1972 \pm 0.1356$ & $0.9825 \pm 0.0004$ \\
PM rat. & 0.05 & 0 & $32.0401 \pm 0.0458$ & $0.6970 \pm 0.0019$ \\
PM rat. & 0.05 & 1 & $35.5576 \pm 0.0616$ & $0.8368 \pm 0.0021$ \\
PM rat. & 0.05 & 2 & $38.8550 \pm 0.0711$ & $0.9213 \pm 0.0016$ \\
PM rat. & 0.05 & 5 & $43.4407 \pm 0.1462$ & $0.9799 \pm 0.0009$ \\
PM rat. & 0.05 & 10 & $43.6020 \pm 0.1602$ & $0.9884 \pm 0.0004$ \\
PM rat. & 0.05 & 20 & $40.4397 \pm 0.1318$ & $0.9834 \pm 0.0004$ \\
PM rat. & 0.05 & 40 & $35.8450 \pm 0.1134$ & $0.9621 \pm 0.0007$ \\
PM rat. & 0.05 & 80 & $31.9488 \pm 0.0883$ & $0.9306 \pm 0.0007$ \\
PM rat. & 0.1 & 0 & $32.0401 \pm 0.0458$ & $0.6970 \pm 0.0019$ \\
PM rat. & 0.1 & 1 & $37.5736 \pm 0.0566$ & $0.8951 \pm 0.0015$ \\
PM rat. & 0.1 & 2 & $40.2716 \pm 0.1021$ & $0.9517 \pm 0.0013$ \\
PM rat. & 0.1 & 5 & $40.0917 \pm 0.1458$ & $0.9763 \pm 0.0009$ \\
PM rat. & 0.1 & 10 & $36.9634 \pm 0.1010$ & $0.9717 \pm 0.0005$ \\
PM rat. & 0.1 & 20 & $33.2771 \pm 0.0674$ & $0.9500 \pm 0.0004$ \\
PM rat. & 0.1 & 40 & $29.5151 \pm 0.0394$ & $0.9113 \pm 0.0005$ \\
PM rat. & 0.1 & 80 & $25.0462 \pm 0.0230$ & $0.8433 \pm 0.0005$ \\
PM rat. & 0.2 & 0 & $32.0401 \pm 0.0458$ & $0.6970 \pm 0.0019$ \\
PM rat. & 0.2 & 1 & $37.6875 \pm 0.0724$ & $0.9115 \pm 0.0015$ \\
PM rat. & 0.2 & 2 & $37.9034 \pm 0.1048$ & $0.9510 \pm 0.0013$ \\
PM rat. & 0.2 & 5 & $34.3099 \pm 0.0614$ & $0.9604 \pm 0.0008$ \\
PM rat. & 0.2 & 10 & $30.3356 \pm 0.0320$ & $0.9390 \pm 0.0004$ \\
PM rat. & 0.2 & 20 & $26.0913 \pm 0.0219$ & $0.8849 \pm 0.0003$ \\
PM rat. & 0.2 & 40 & $23.1674 \pm 0.0115$ & $0.8145 \pm 0.0003$ \\
PM rat. & 0.2 & 80 & $21.3571 \pm 0.0068$ & $0.7574 \pm 0.0003$ \\
\end{longtable}

\begin{longtable}{llrll}
\caption{Rampe et contours, $\sigma=0.05$ : grille complète, moyennes $\pm$ écarts-types sur dix graines.}\\
\toprule
Méthode & $K$ & $n$ & PSNR (dB) & SSIM \\
\midrule
\endfirsthead
\toprule
Méthode & $K$ & $n$ & PSNR (dB) & SSIM \\
\midrule
\endhead
\bottomrule
\endfoot
Chaleur & -- & 0 & $26.0195 \pm 0.0458$ & $0.4242 \pm 0.0017$ \\
Chaleur & -- & 1 & $28.2815 \pm 0.0409$ & $0.7223 \pm 0.0039$ \\
Chaleur & -- & 2 & $28.7674 \pm 0.0327$ & $0.8286 \pm 0.0038$ \\
Chaleur & -- & 5 & $27.0963 \pm 0.0231$ & $0.8837 \pm 0.0025$ \\
Chaleur & -- & 10 & $25.6111 \pm 0.0193$ & $0.8782 \pm 0.0010$ \\
Chaleur & -- & 20 & $23.9361 \pm 0.0161$ & $0.8421 \pm 0.0005$ \\
Chaleur & -- & 40 & $22.3730 \pm 0.0139$ & $0.7918 \pm 0.0006$ \\
Chaleur & -- & 80 & $21.0465 \pm 0.0115$ & $0.7483 \pm 0.0006$ \\
Bruitée & -- & 0 & $26.0195 \pm 0.0458$ & $0.4242 \pm 0.0017$ \\
PM exp. & 0.025 & 0 & $26.0195 \pm 0.0458$ & $0.4242 \pm 0.0017$ \\
PM exp. & 0.025 & 1 & $26.1085 \pm 0.0471$ & $0.4282 \pm 0.0018$ \\
PM exp. & 0.025 & 2 & $26.1913 \pm 0.0483$ & $0.4321 \pm 0.0019$ \\
PM exp. & 0.025 & 5 & $26.4087 \pm 0.0511$ & $0.4423 \pm 0.0020$ \\
PM exp. & 0.025 & 10 & $26.7042 \pm 0.0559$ & $0.4567 \pm 0.0023$ \\
PM exp. & 0.025 & 20 & $27.1635 \pm 0.0666$ & $0.4804 \pm 0.0029$ \\
PM exp. & 0.025 & 40 & $27.8380 \pm 0.0834$ & $0.5173 \pm 0.0040$ \\
PM exp. & 0.025 & 80 & $28.6568 \pm 0.0866$ & $0.5631 \pm 0.0043$ \\
PM exp. & 0.05 & 0 & $26.0195 \pm 0.0458$ & $0.4242 \pm 0.0017$ \\
PM exp. & 0.05 & 1 & $26.5860 \pm 0.0533$ & $0.4505 \pm 0.0021$ \\
PM exp. & 0.05 & 2 & $27.1559 \pm 0.0613$ & $0.4784 \pm 0.0026$ \\
PM exp. & 0.05 & 5 & $28.8571 \pm 0.0880$ & $0.5692 \pm 0.0044$ \\
PM exp. & 0.05 & 10 & $31.3952 \pm 0.1285$ & $0.7084 \pm 0.0066$ \\
PM exp. & 0.05 & 20 & $34.2341 \pm 0.1965$ & $0.8345 \pm 0.0077$ \\
PM exp. & 0.05 & 40 & $35.9884 \pm 0.2454$ & $0.8901 \pm 0.0081$ \\
PM exp. & 0.05 & 80 & $36.8297 \pm 0.2473$ & $0.9165 \pm 0.0068$ \\
PM exp. & 0.1 & 0 & $26.0195 \pm 0.0458$ & $0.4242 \pm 0.0017$ \\
PM exp. & 0.1 & 1 & $28.4276 \pm 0.0653$ & $0.5408 \pm 0.0031$ \\
PM exp. & 0.1 & 2 & $30.9744 \pm 0.0855$ & $0.6740 \pm 0.0045$ \\
PM exp. & 0.1 & 5 & $36.7022 \pm 0.1246$ & $0.9028 \pm 0.0035$ \\
PM exp. & 0.1 & 10 & $39.2732 \pm 0.1910$ & $0.9632 \pm 0.0019$ \\
PM exp. & 0.1 & 20 & $38.1642 \pm 0.2302$ & $0.9699 \pm 0.0014$ \\
PM exp. & 0.1 & 40 & $35.8410 \pm 0.2385$ & $0.9580 \pm 0.0014$ \\
PM exp. & 0.1 & 80 & $33.7561 \pm 0.2474$ & $0.9436 \pm 0.0020$ \\
PM exp. & 0.2 & 0 & $26.0195 \pm 0.0458$ & $0.4242 \pm 0.0017$ \\
PM exp. & 0.2 & 1 & $31.2776 \pm 0.0638$ & $0.6842 \pm 0.0034$ \\
PM exp. & 0.2 & 2 & $34.5041 \pm 0.0993$ & $0.8342 \pm 0.0036$ \\
PM exp. & 0.2 & 5 & $36.3603 \pm 0.1915$ & $0.9304 \pm 0.0029$ \\
PM exp. & 0.2 & 10 & $35.3247 \pm 0.2146$ & $0.9501 \pm 0.0014$ \\
PM exp. & 0.2 & 20 & $33.4172 \pm 0.1385$ & $0.9446 \pm 0.0007$ \\
PM exp. & 0.2 & 40 & $31.3213 \pm 0.1233$ & $0.9266 \pm 0.0011$ \\
PM exp. & 0.2 & 80 & $28.4888 \pm 0.2452$ & $0.8948 \pm 0.0042$ \\
PM rat. & 0.025 & 0 & $26.0195 \pm 0.0458$ & $0.4242 \pm 0.0017$ \\
PM rat. & 0.025 & 1 & $26.5708 \pm 0.0503$ & $0.4491 \pm 0.0020$ \\
PM rat. & 0.025 & 2 & $27.1429 \pm 0.0550$ & $0.4763 \pm 0.0023$ \\
PM rat. & 0.025 & 5 & $28.9892 \pm 0.0708$ & $0.5704 \pm 0.0035$ \\
PM rat. & 0.025 & 10 & $32.4470 \pm 0.1049$ & $0.7479 \pm 0.0053$ \\
PM rat. & 0.025 & 20 & $38.8334 \pm 0.1985$ & $0.9482 \pm 0.0027$ \\
PM rat. & 0.025 & 40 & $42.0671 \pm 0.2953$ & $0.9874 \pm 0.0009$ \\
PM rat. & 0.025 & 80 & $39.5634 \pm 0.2819$ & $0.9810 \pm 0.0008$ \\
PM rat. & 0.05 & 0 & $26.0195 \pm 0.0458$ & $0.4242 \pm 0.0017$ \\
PM rat. & 0.05 & 1 & $27.5840 \pm 0.0560$ & $0.4973 \pm 0.0025$ \\
PM rat. & 0.05 & 2 & $29.2693 \pm 0.0667$ & $0.5833 \pm 0.0033$ \\
PM rat. & 0.05 & 5 & $34.4403 \pm 0.1011$ & $0.8295 \pm 0.0043$ \\
PM rat. & 0.05 & 10 & $39.0370 \pm 0.1915$ & $0.9561 \pm 0.0020$ \\
PM rat. & 0.05 & 20 & $39.0639 \pm 0.2486$ & $0.9754 \pm 0.0011$ \\
PM rat. & 0.05 & 40 & $35.5615 \pm 0.2303$ & $0.9597 \pm 0.0015$ \\
PM rat. & 0.05 & 80 & $31.8722 \pm 0.1823$ & $0.9298 \pm 0.0015$ \\
PM rat. & 0.1 & 0 & $26.0195 \pm 0.0458$ & $0.4242 \pm 0.0017$ \\
PM rat. & 0.1 & 1 & $29.5405 \pm 0.0604$ & $0.5957 \pm 0.0031$ \\
PM rat. & 0.1 & 2 & $32.8048 \pm 0.0740$ & $0.7599 \pm 0.0038$ \\
PM rat. & 0.1 & 5 & $36.7886 \pm 0.1771$ & $0.9249 \pm 0.0031$ \\
PM rat. & 0.1 & 10 & $36.0280 \pm 0.1819$ & $0.9541 \pm 0.0016$ \\
PM rat. & 0.1 & 20 & $33.0595 \pm 0.1337$ & $0.9447 \pm 0.0008$ \\
PM rat. & 0.1 & 40 & $29.4453 \pm 0.0775$ & $0.9096 \pm 0.0010$ \\
PM rat. & 0.1 & 80 & $25.0265 \pm 0.0454$ & $0.8426 \pm 0.0009$ \\
PM rat. & 0.2 & 0 & $26.0195 \pm 0.0458$ & $0.4242 \pm 0.0017$ \\
PM rat. & 0.2 & 1 & $31.5236 \pm 0.0613$ & $0.7002 \pm 0.0032$ \\
PM rat. & 0.2 & 2 & $34.0080 \pm 0.1089$ & $0.8383 \pm 0.0036$ \\
PM rat. & 0.2 & 5 & $33.3759 \pm 0.1154$ & $0.9212 \pm 0.0027$ \\
PM rat. & 0.2 & 10 & $30.1195 \pm 0.0668$ & $0.9242 \pm 0.0012$ \\
PM rat. & 0.2 & 20 & $26.0485 \pm 0.0443$ & $0.8801 \pm 0.0006$ \\
PM rat. & 0.2 & 40 & $23.1546 \pm 0.0230$ & $0.8130 \pm 0.0006$ \\
PM rat. & 0.2 & 80 & $21.3519 \pm 0.0137$ & $0.7568 \pm 0.0006$ \\
\end{longtable}

\begin{longtable}{llrll}
\caption{Rampe et contours, $\sigma=0.1$ : grille complète, moyennes $\pm$ écarts-types sur dix graines.}\\
\toprule
Méthode & $K$ & $n$ & PSNR (dB) & SSIM \\
\midrule
\endfirsthead
\toprule
Méthode & $K$ & $n$ & PSNR (dB) & SSIM \\
\midrule
\endhead
\bottomrule
\endfoot
Chaleur & -- & 0 & $19.9989 \pm 0.0458$ & $0.2314 \pm 0.0011$ \\
Chaleur & -- & 1 & $25.2198 \pm 0.0662$ & $0.4564 \pm 0.0043$ \\
Chaleur & -- & 2 & $26.8227 \pm 0.0707$ & $0.6136 \pm 0.0064$ \\
Chaleur & -- & 5 & $26.4698 \pm 0.0548$ & $0.7762 \pm 0.0066$ \\
Chaleur & -- & 10 & $25.3685 \pm 0.0409$ & $0.8298 \pm 0.0033$ \\
Chaleur & -- & 20 & $23.8463 \pm 0.0323$ & $0.8247 \pm 0.0015$ \\
Chaleur & -- & 40 & $22.3370 \pm 0.0277$ & $0.7862 \pm 0.0014$ \\
Chaleur & -- & 80 & $21.0311 \pm 0.0232$ & $0.7464 \pm 0.0012$ \\
Bruitée & -- & 0 & $19.9989 \pm 0.0458$ & $0.2314 \pm 0.0011$ \\
PM exp. & 0.025 & 0 & $19.9989 \pm 0.0458$ & $0.2314 \pm 0.0011$ \\
PM exp. & 0.025 & 1 & $20.0108 \pm 0.0459$ & $0.2317 \pm 0.0011$ \\
PM exp. & 0.025 & 2 & $20.0217 \pm 0.0461$ & $0.2320 \pm 0.0011$ \\
PM exp. & 0.025 & 5 & $20.0489 \pm 0.0465$ & $0.2327 \pm 0.0011$ \\
PM exp. & 0.025 & 10 & $20.0819 \pm 0.0471$ & $0.2335 \pm 0.0011$ \\
PM exp. & 0.025 & 20 & $20.1256 \pm 0.0480$ & $0.2346 \pm 0.0011$ \\
PM exp. & 0.025 & 40 & $20.1794 \pm 0.0491$ & $0.2360 \pm 0.0012$ \\
PM exp. & 0.025 & 80 & $20.2439 \pm 0.0504$ & $0.2377 \pm 0.0012$ \\
PM exp. & 0.05 & 0 & $19.9989 \pm 0.0458$ & $0.2314 \pm 0.0011$ \\
PM exp. & 0.05 & 1 & $20.0881 \pm 0.0469$ & $0.2337 \pm 0.0011$ \\
PM exp. & 0.05 & 2 & $20.1712 \pm 0.0479$ & $0.2358 \pm 0.0011$ \\
PM exp. & 0.05 & 5 & $20.3894 \pm 0.0503$ & $0.2415 \pm 0.0012$ \\
PM exp. & 0.05 & 10 & $20.6864 \pm 0.0540$ & $0.2497 \pm 0.0013$ \\
PM exp. & 0.05 & 20 & $21.1501 \pm 0.0622$ & $0.2635 \pm 0.0015$ \\
PM exp. & 0.05 & 40 & $21.8340 \pm 0.0788$ & $0.2869 \pm 0.0023$ \\
PM exp. & 0.05 & 80 & $22.6527 \pm 0.0862$ & $0.3189 \pm 0.0035$ \\
PM exp. & 0.1 & 0 & $19.9989 \pm 0.0458$ & $0.2314 \pm 0.0011$ \\
PM exp. & 0.1 & 1 & $20.5670 \pm 0.0529$ & $0.2461 \pm 0.0013$ \\
PM exp. & 0.1 & 2 & $21.1386 \pm 0.0603$ & $0.2624 \pm 0.0015$ \\
PM exp. & 0.1 & 5 & $22.8441 \pm 0.0852$ & $0.3224 \pm 0.0031$ \\
PM exp. & 0.1 & 10 & $25.3782 \pm 0.1243$ & $0.4540 \pm 0.0079$ \\
PM exp. & 0.1 & 20 & $28.0985 \pm 0.1844$ & $0.6437 \pm 0.0131$ \\
PM exp. & 0.1 & 40 & $29.4064 \pm 0.2339$ & $0.7469 \pm 0.0167$ \\
PM exp. & 0.1 & 80 & $29.6056 \pm 0.2034$ & $0.7946 \pm 0.0120$ \\
PM exp. & 0.2 & 0 & $19.9989 \pm 0.0458$ & $0.2314 \pm 0.0011$ \\
PM exp. & 0.2 & 1 & $22.4143 \pm 0.0652$ & $0.3011 \pm 0.0019$ \\
PM exp. & 0.2 & 2 & $24.9662 \pm 0.0866$ & $0.4054 \pm 0.0039$ \\
PM exp. & 0.2 & 5 & $30.5802 \pm 0.1171$ & $0.7254 \pm 0.0074$ \\
PM exp. & 0.2 & 10 & $32.8772 \pm 0.2265$ & $0.8781 \pm 0.0044$ \\
PM exp. & 0.2 & 20 & $32.3930 \pm 0.2530$ & $0.9209 \pm 0.0023$ \\
PM exp. & 0.2 & 40 & $30.5046 \pm 0.3715$ & $0.9143 \pm 0.0044$ \\
PM exp. & 0.2 & 80 & $27.9716 \pm 0.3660$ & $0.8860 \pm 0.0069$ \\
PM rat. & 0.025 & 0 & $19.9989 \pm 0.0458$ & $0.2314 \pm 0.0011$ \\
PM rat. & 0.025 & 1 & $20.1674 \pm 0.0472$ & $0.2356 \pm 0.0011$ \\
PM rat. & 0.025 & 2 & $20.3370 \pm 0.0486$ & $0.2398 \pm 0.0011$ \\
PM rat. & 0.025 & 5 & $20.8536 \pm 0.0529$ & $0.2535 \pm 0.0013$ \\
PM rat. & 0.025 & 10 & $21.7496 \pm 0.0606$ & $0.2798 \pm 0.0016$ \\
PM rat. & 0.025 & 20 & $23.7056 \pm 0.0794$ & $0.3522 \pm 0.0031$ \\
PM rat. & 0.025 & 40 & $28.1606 \pm 0.1311$ & $0.6015 \pm 0.0094$ \\
PM rat. & 0.025 & 80 & $34.5137 \pm 0.3117$ & $0.9229 \pm 0.0053$ \\
PM rat. & 0.05 & 0 & $19.9989 \pm 0.0458$ & $0.2314 \pm 0.0011$ \\
PM rat. & 0.05 & 1 & $20.5522 \pm 0.0500$ & $0.2453 \pm 0.0012$ \\
PM rat. & 0.05 & 2 & $21.1263 \pm 0.0544$ & $0.2609 \pm 0.0013$ \\
PM rat. & 0.05 & 5 & $22.9795 \pm 0.0689$ & $0.3214 \pm 0.0022$ \\
PM rat. & 0.05 & 10 & $26.4432 \pm 0.1027$ & $0.4864 \pm 0.0064$ \\
PM rat. & 0.05 & 20 & $32.5651 \pm 0.1841$ & $0.8422 \pm 0.0070$ \\
PM rat. & 0.05 & 40 & $34.2340 \pm 0.2965$ & $0.9465 \pm 0.0025$ \\
PM rat. & 0.05 & 80 & $31.4959 \pm 0.2962$ & $0.9265 \pm 0.0027$ \\
PM rat. & 0.1 & 0 & $19.9989 \pm 0.0458$ & $0.2314 \pm 0.0011$ \\
PM rat. & 0.1 & 1 & $21.5693 \pm 0.0556$ & $0.2733 \pm 0.0014$ \\
PM rat. & 0.1 & 2 & $23.2604 \pm 0.0662$ & $0.3300 \pm 0.0022$ \\
PM rat. & 0.1 & 5 & $28.4109 \pm 0.1065$ & $0.5911 \pm 0.0070$ \\
PM rat. & 0.1 & 10 & $32.5660 \pm 0.2041$ & $0.8535 \pm 0.0056$ \\
PM rat. & 0.1 & 20 & $32.1296 \pm 0.2092$ & $0.9185 \pm 0.0023$ \\
PM rat. & 0.1 & 40 & $29.1451 \pm 0.1527$ & $0.9018 \pm 0.0022$ \\
PM rat. & 0.1 & 80 & $24.9353 \pm 0.0958$ & $0.8396 \pm 0.0021$ \\
PM rat. & 0.2 & 0 & $19.9989 \pm 0.0458$ & $0.2314 \pm 0.0011$ \\
PM rat. & 0.2 & 1 & $23.5328 \pm 0.0612$ & $0.3384 \pm 0.0020$ \\
PM rat. & 0.2 & 2 & $26.7875 \pm 0.0802$ & $0.4911 \pm 0.0044$ \\
PM rat. & 0.2 & 5 & $30.4631 \pm 0.1692$ & $0.7727 \pm 0.0072$ \\
PM rat. & 0.2 & 10 & $29.2988 \pm 0.1374$ & $0.8644 \pm 0.0040$ \\
PM rat. & 0.2 & 20 & $25.8938 \pm 0.0902$ & $0.8609 \pm 0.0016$ \\
PM rat. & 0.2 & 40 & $23.1136 \pm 0.0462$ & $0.8073 \pm 0.0013$ \\
PM rat. & 0.2 & 80 & $21.3363 \pm 0.0279$ & $0.7550 \pm 0.0012$ \\
\end{longtable}


